% ===== UNIT 34: What drives the ViT-B gain (6.1) =====
% Status: AI-DRAFT  (change to REWRITTEN / REVIEWED when done; see CHECKLIST.md)
% Rules: style rewrite only. Keep every \cite, \ref, \label, number, equation and table value as is. Change the wording, not the content.
\subsection{What drives the ViT-B gain}
The improvement of our ViT-B pipeline over the published TransReID baseline
combines recipe and post-processing effects that the ablation deliberately
separates. Within the recipe, CLIP pretraining initialisation and AdamW with
layer-wise LR decay are the two ingredients our development depended on most
(A2, A4); CLIP provides a rich visual-semantic prior for fine-grained vehicle
discrimination, consistent with transfer-learning observations in the ReID
literature~\cite{li2023clipreid,radford2021clip}. We reiterate the caveat of
Section~\ref{sec:matched}: at matched no-post-processing conditions the recipe
does not by itself exceed TransReID (82.2\% vs.\ 82.3\% mAP), and under our
identical post-processing both TransReID and CLIP-ReID in fact surpass our
pipeline (91.12\% and 91.80\% vs.\ 89.97\%). These are therefore drivers
within our recipe rather than net gains over a tuned baseline, and the
bulk of the reported improvement is contributed by AQE, $k$-reciprocal
re-ranking, and TTA. For context, a DINOv2-initialised
alternative~\cite{oquab2023dinov2} evaluated as a second fusion stream on
VeRi-776 did not match the CLIP-based CLIP-SENet stream (below 89.3\% mAP
standalone), so CLIP-family semantic supervision was retained for both streams;
we did not run a controlled DINOv2-versus-CLIP Stream-1 backbone swap on
VeRi-776 and therefore make no claim about pretraining scale versus semantic
supervision from this benchmark alone.
